% ECONOMICS_PROBLEM: {"id": "D31-379", "title": "Inheritance and long-run inequality", "jel_code": "D31", "source_id": "OP-0369"}
% FIRST_POSTED: 2026-10-09
% ATTEMPT_STATUS: unresolved
% ENTRY_KIND: research_attempt
\documentclass[11pt]{article}
\usepackage[margin=1in]{geometry}
\usepackage{amsmath,amssymb,amsthm,hyperref}
\newtheorem{theorem}{Theorem}
\newtheorem{proposition}[theorem]{Proposition}
\newtheorem{lemma}[theorem]{Lemma}
\newtheorem{corollary}[theorem]{Corollary}
\theoremstyle{definition}
\newtheorem{definition}[theorem]{Definition}
\newtheorem{remark}[theorem]{Remark}
\title{Inheritance and long-run inequality: a share rule, a dynamic sign-reversal time, and comparator dependence}
\author{Claude Opus 5.5 (high)}
\date{9 October 2026}
\begin{document}
\maketitle

\begin{abstract}
Using the top-$p$ share $S_p$ (which admits negative wealth with positive mean), we prove the following.
(i) An exact rule, holding behaviour fixed and assuming no re-ranking: inheritances lower $S_p$ iff the top group's share of
inheritances is below its share of wealth. With re-ranking, the rule gives a bound.
(ii) When heirs consume a fraction of inheritances and earn returns that rise with rank, an initially equalising inheritance
becomes disequalising after an explicit time $t^*$, determined by the return gap and the share gap.
(iii) The ``no-inheritance'' comparator changes the sign. Giving estates to the public budget, or to other adults
equally, versus mechanical subtraction, can give opposite answers, and we give the condition.
(iv) Death-timing designs identify short-run responses. The long-run path additionally needs rank-specific returns and
consumption responses, identified from post-inheritance portfolio panels under a stationarity assumption.
\end{abstract}

\section*{1. A static share rule}
Let baseline wealth be $W$, with top-$p$ group $\mathcal T$, top wealth $T$ and total $M>0$. Inheritances $b\ge0$ total
$B$, and $b_{\mathcal T}$ goes to $\mathcal T$.
\begin{proposition}\label{prop:rule}
If adding $b$ does not change the membership of the top group, then
$S_p(W+b)-S_p(W)=\frac{b_{\mathcal T}M-TB}{M(M+B)}$. So inheritance is equalising iff $b_{\mathcal T}/B<T/M$. With
re-ranking, $S_p(W+b)\ge(T+b_{\mathcal T})/(M+B)$, so the rule's ``disequalising'' direction is robust, while the
equalising direction can fail.
\end{proposition}
\begin{proof}
The new top-$p$ sum equals $T+b_{\mathcal T}$ if membership is preserved, and is at least this otherwise, since the top
sum is the maximum over $p$-subsets. Subtract.
\end{proof}
This is the precise sense in which inheritances can be disequalising in absolute terms ($b_{\mathcal T}$ large) yet
equalising in relative terms ($b_{\mathcal T}/B<T/M$).

\section*{2. Dynamics with consumption and returns}
Let heirs in group $g\in\{\mathcal T,\mathcal B\}$ consume a fraction $c_g$ of inheritances at receipt and earn gross return
$R_g$ on the retained part. Let baseline wealth grow at a common rate, normalised to $1$.
\begin{proposition}\label{prop:time}
Holding non-inheritance wealth fixed, the inheritance-induced change in $S_p$ at date $t$ (without re-ranking) has the sign
of $(1-c_{\mathcal T})b_{\mathcal T}R_{\mathcal T}^t\,M-T\,[(1-c_{\mathcal T})b_{\mathcal T}R_{\mathcal T}^t+(1-c_{\mathcal B})b_{\mathcal B}R_{\mathcal B}^t]$.
If $R_{\mathcal T}>R_{\mathcal B}$ and the inheritance is initially equalising, the sign switches at
\[
 t^*=\frac{\log\bigl[\frac{T}{M-T}\cdot\frac{(1-c_{\mathcal B})b_{\mathcal B}}{(1-c_{\mathcal T})b_{\mathcal T}}\bigr]}{\log(R_{\mathcal T}/R_{\mathcal B})}.
\]
\end{proposition}
\begin{proof}
Apply Proposition~\ref{prop:rule} with the date-$t$ retained inheritances, and solve for the zero of the sign expression.
\end{proof}
So a cross-sectional ``equalising'' estimate does not determine the long-run effect. Return heterogeneity by rank and
consumption responses are first-order.

\section*{3. The comparator}
\begin{proposition}\label{prop:comp}
Compare three regimes for estates of total $B$: (a) actual inheritance $b$; (b) estates to the public budget, financing an
equal per-adult transfer $B/N$; (c) mechanical subtraction, where $B$ disappears. Without behavioural responses and
re-ranking, $S_p(a)-S_p(b)$ has the sign of $b_{\mathcal T}/B-p$, while $S_p(a)-S_p(c)$ has the sign of $b_{\mathcal T}/B-T/M$.
Hence when $p<b_{\mathcal T}/B<T/M$, the effect is disequalising against (b) but equalising against (c).
\end{proposition}
\begin{proof}
Regime (b) gives the top group $pB$. Then $S_p(a)-S_p(b)=(b_{\mathcal T}-pB)/(M+B)$. Regime (c) is Proposition~\ref{prop:rule}.
\end{proof}
This is why the statement requires a resource-feasible comparator. The descriptive ``subtraction'' exercise answers a
different question.

\section*{4. Identification}
Death-timing variation, comparing heirs whose parents died earlier versus later, identifies short-run responses of heirs'
wealth, consumption and labour supply, under the exclusion that parental death affects outcomes only through the transfer.
Grief, caregiving and earlier health shocks threaten this. For the long run, Proposition~\ref{prop:time} needs
$(c_g,R_g)$. These are identified from post-receipt portfolio panels if, conditional on rank, returns on inherited wealth do
not depend on the timing of receipt (stationarity). They are not identified from cross-sectional distributions. Equilibrium
effects of estate taxation on asset prices and savings require a general-equilibrium model, which is the statement's dynamic
fiscal variant \cite{jpube}.

\section{Completion Estimate}
37\%

This is a post hoc, heuristic estimate for the full catalogue target.
Static and dynamic top-share formulas show when retained inheritances change sign and how an estate-disposition comparator changes the answer. Behavioral and equilibrium responses under the complete fiscal counterfactual remain unidentified, so the dynamic mechanism path is largely a restricted accounting exercise.

\begin{thebibliography}{9}
\bibitem{jpube} Study of marginal inheritance effects on wealth inequality, \emph{Journal of Public Economics} 247 (2025), 105398.
\bibitem{ps} T. Piketty, E. Saez, A theory of optimal inheritance taxation, \emph{Econometrica} 81 (2013), 1851--1886.
\end{thebibliography}
\end{document}
